\documentclass[11pt,a4paper]{article}
\usepackage[margin=2.5cm]{geometry}
\usepackage{amsmath,amsthm}
\usepackage{newtxtext}
\usepackage{newtxmath}
\usepackage{enumitem}
\usepackage{microtype}
\setlength{\parindent}{0pt}
\setlength{\parskip}{4pt}
\allowdisplaybreaks

\newcommand{\E}{\mathbb{E}}
\renewcommand{\P}{\mathbb{P}}
\newcommand{\R}{\mathbb{R}}
\newcommand{\N}{\mathbb{N}}
\newcommand{\Q}{\mathbb{Q}}
\newcommand{\F}{\mathcal{F}}
\newcommand{\G}{\mathcal{G}}
\newcommand{\Hc}{\mathcal{H}}
\newcommand{\1}{\mathbf{1}}
\newcommand{\Leb}{\mathrm{Leb}}
\newcommand{\sgn}{\operatorname{sgn}}
\newcommand{\prob}[1]{\par\medskip\noindent\textbf{#1}\quad}
\newcommand{\sub}[1]{\par\smallskip\noindent\textbf{#1}\ }
\newcommand{\note}[1]{\par\noindent\emph{Remark.} #1}

\begin{document}

\begin{center}
{\large\textbf{Stochastic Processes --- 1060429}}\\[2pt]
Nawar Sgier, 326586047
\end{center}

\noindent Throughout, a Brownian motion (BM) is standard ($B_0=0$) with continuous paths, and
$\F_t$ denotes the filtration in use. Apparent typos are pointed out where they occur.

\section*{Exercise 4}
\subsection*{Stochastic calculus}

Throughout, for $f\in L^2$ (adapted, with $\E\int_0^tf^2ds<\infty$), $\int_0^tf\,dB$ is the $L^2$-limit of
$\int f_n\,dB$ for simple adapted $f_n\to f$ in $L^2(ds\otimes d\P)$. For a partition
$0=t_0<\dots<t_n=t$ of mesh $|\pi|$ write $\Delta_i=B_{t_i}-B_{t_{i-1}}$ and $\Delta t_i=t_i-t_{i-1}$.

\prob{1.}
\note{In (a)ii and (b)i the right-hand side should be $\frac13B_t^3-\int_0^tB_s\,ds$, not $\frac13B_t^2-\dots$.
As stated it cannot hold: the left side has mean $0$, while $\E[\frac13B_t^2-\int_0^tB_s\,ds]=\frac t3$.
Itô's formula in (b) confirms the corrected form.}

\sub{(a) i.} $f_\pi=\sum t_{i-1}\1_{[t_{i-1},t_i)}\to s$ in $L^2$, since $\int_0^t|f_\pi(s)-s|^2ds\le t|\pi|^2$. Summation by parts gives
\[
\sum_it_{i-1}\Delta_i=tB_t-\sum_iB_{t_i}\Delta t_i .
\]
By continuity $|\sum_iB_{t_i}\Delta t_i-\int_0^tB_s\,ds|\to0$ pathwise. It is bounded by $2t\sup_{[0,t]}|B|\in L^2$,
so the convergence also holds in $L^2$. Hence $\int_0^ts\,dB_s=tB_t-\int_0^tB_s\,ds$.

\sub{(a) ii.} For $u<s\le t$, by Cauchy--Schwarz and $B_s+B_u\sim N(0,s+3u)$,
$\E(B_s^2-B_u^2)^2\le\sqrt{3(s-u)^2}\sqrt{3(4t)^2}=12t(s-u)$. Hence
$\sum_iB_{t_{i-1}}^2\1_{[t_{i-1},t_i)}\to B^2$ in $L^2(ds\otimes d\P)$, and $\int_0^tB^2dB=\lim_{L^2}S_\pi$
with $S_\pi=\sum_iB_{t_{i-1}}^2\Delta_i$. Summing $b^3-a^3=3a^2(b-a)+3a(b-a)^2+(b-a)^3$ over the partition gives
\[
S_\pi=\tfrac13B_t^3-\sum_iB_{t_{i-1}}\Delta_i^2-\tfrac13\sum_i\Delta_i^3 .
\]
The cross terms below vanish by independence of increments and since odd moments are $0$, resp.\ $\E[\Delta_j^2-\Delta t_j]=0$:
\[
\E\Big(\sum_i\Delta_i^3\Big)^2=\sum_i15\Delta t_i^3\le15t|\pi|^2,\qquad
\E\Big(\sum_iB_{t_{i-1}}(\Delta_i^2-\Delta t_i)\Big)^2=\sum_it_{i-1}\cdot2\Delta t_i^2\le2t^2|\pi| .
\]
Moreover $\sum_iB_{t_{i-1}}\Delta t_i\to\int_0^tB_s\,ds$ in $L^2$ as in i. Hence
$\int_0^tB_s^2\,dB_s=\frac13B_t^3-\int_0^tB_s\,ds$.

\sub{(b) i.} Itô's formula with $f(x)=x^3/3$ (so $f'=x^2$, $f''=2x$) gives
$\frac13B_t^3=\int_0^tB_s^2\,dB_s+\int_0^tB_s\,ds$.

\prob{2.}
\note{We take $f,g\in L^2([0,T]\times\Omega)$, i.e.\ $\E\int_0^Tf^2<\infty$ for every $T$, the class on which the Itô integral is
an $L^2$-martingale. If only $\int_0^Tf^2ds<\infty$ a.s.\ were assumed, the claim would be false in
general. By Dudley's representation theorem (not part of the course and not used here) there is such an $f$ with
$\int_0^Tf\,dB=1$ a.s., and then $0+\int_0^Tf\,dB=1+\int_0^T0\,dB$ with $C\ne D$.}

The Itô integral has mean zero, so taking expectations gives $C=D$. Then $\int_0^T(f-g)\,dB=0$ a.s., and by the Itô
isometry
\[\E\int_0^T(f-g)^2ds=\E\Big(\int_0^T(f-g)\,dB\Big)^2=0,\]
i.e.\ $f=g$ for $ds\otimes d\P$-a.e.\ $(s,\omega)$.

\prob{3.}
\sub{(a)} $d(B_t^2)=dt+2B_t\,dB_t$.

\sub{(b)} $d(2+t+e^{B_t})=(1+\tfrac12e^{B_t})\,dt+e^{B_t}\,dB_t$.

\sub{(c)} $d(B_t^2+\tilde B_t^2)=2\,dt+2B_t\,dB_t+2\tilde B_t\,d\tilde B_t$, i.e.\ $\alpha=2$ and $\beta=(2B_t,2\tilde B_t)$ against
$d(B_t,\tilde B_t)$.

\sub{(d)} $d(t,B_t)=\binom10dt+\binom01dB_t$.

\sub{(e)} With $f(t,x)=e^{t/2}\cos x$: $f_t=\frac12f$, $f_x=-e^{t/2}\sin x$, $f_{xx}=-f$. Hence
\[
d\big(e^{t/2}\cos B_t\big)=\big(\tfrac12f-\tfrac12f\big)dt-e^{t/2}\sin B_t\,dB_t=-e^{t/2}\sin B_t\,dB_t .
\]
So $e^{t/2}\cos B_t=1-\int_0^te^{s/2}\sin B_s\,dB_s$. The integrand is bounded by $e^{T/2}$ on $[0,T]$, so it lies in $L^2$ and
the Itô integral is a martingale. Therefore $e^{t/2}\cos B_t$ is a martingale.

\prob{4.}
\sub{(a)} With $f=e^x$, $f'=f''=f$: $dX_t=e^{B_t}dB_t+\frac12e^{B_t}dt=\frac12X_tdt+X_tdB_t$.

\sub{(b)} With $f=\sin x$: $dX_t=\cos B_t\,dB_t-\frac12\sin B_t\,dt$. While $\cos B_t>0$, i.e.\ for
$t<\inf\{s:|B_s|=\pi/2\}$, we have $\cos B_t=\sqrt{1-\sin^2B_t}=\sqrt{1-X_t^2}$. Hence
$dX_t=-\frac12X_tdt+\sqrt{1-X_t^2}\,dB_t$.

\prob{5.}
For $k\ge2$, Itô's formula gives
\[B_t^k=\int_0^tkB_s^{k-1}dB_s+\tfrac12k(k-1)\int_0^tB_s^{k-2}ds.\]
The integrand satisfies $\E\int_0^tB_s^{2k-2}ds<\infty$, so the stochastic integral has mean $0$. By Fubini,
\[\beta_k(t)=\tfrac12k(k-1)\int_0^t\beta_{k-2}(s)\,ds.\]
The formula also holds for $k=1$ ($\beta_1=0$) but fails for $k=0$, where $\beta_0=1$. With $\beta_0=1$:
$\beta_2(t)=t$, then $\E B_t^4=\beta_4(t)=6\int_0^ts\,ds=3t^2$ and $\E B_t^6=\beta_6(t)=15\int_0^t3s^2\,ds=15t^3$.

\prob{6.}
\note{$g_\varepsilon$ is not $C^2$ on $\R$. It is $C^1$, with
$g_\varepsilon'(x)=\sgn x$ for $|x|\ge\varepsilon$ and $x/\varepsilon$ for $|x|<\varepsilon$. Its second derivative
$g_\varepsilon''=\varepsilon^{-1}\1_{(-\varepsilon,\varepsilon)}$ exists except at $\pm\varepsilon$, so $g_\varepsilon$ is
$C^2$ on $\R\setminus\{\pm\varepsilon\}$. Itô's formula still applies, as shown in (a).}

\sub{(a)} Let $\phi_n$ be continuous, $0\le\phi_n\le1/\varepsilon$, with $\phi_n=g_\varepsilon''$ outside
$(\pm\varepsilon-\frac1n,\pm\varepsilon+\frac1n)$. Let $h_n(x)=g_\varepsilon(0)+\int_0^x\int_0^y\phi_n$, which is $C^2$. Since
$g_\varepsilon'(0)=0$, $|h_n'-g_\varepsilon'|\le\int|\phi_n-g_\varepsilon''|\le4/(n\varepsilon)$ and
$|h_n-g_\varepsilon|(x)\le4|x|/(n\varepsilon)$. Itô's formula for $h_n$ reads
\[
h_n(B_t)=h_n(B_0)+\int_0^th_n'(B_s)\,dB_s+\tfrac12\int_0^t\phi_n(B_s)\,ds .
\]
As $n\to\infty$: $h_n(B_t)\to g_\varepsilon(B_t)$, and $\int h_n'(B)dB\to\int g_\varepsilon'(B)dB$ in $L^2$ by the isometry and
uniform convergence. Also $\phi_n(B_s)\to g_\varepsilon''(B_s)$ for all $s$ with $B_s\ne\pm\varepsilon$, which is a.e.\ $s$ a.s.\
because $\E\,\Leb\{s\le t:B_s=\pm\varepsilon\}=\int_0^t\P(B_s=\pm\varepsilon)ds=0$. Pass to a subsequence along which the stochastic integrals converge a.s.; bounded
convergence for the $ds$-integrals then gives
\[
g_\varepsilon(B_t)=g_\varepsilon(B_0)+\int_0^tg_\varepsilon'(B_s)\,dB_s+\frac1{2\varepsilon}\Leb\{s\in[0,t]:B_s\in(-\varepsilon,\varepsilon)\}.
\]

\sub{(b)} By the isometry and dominated convergence ($\P(B_s=0)=0$ for $s>0$),
\[
\E\Big(\int_0^t\tfrac{B_s}{\varepsilon}\1_{\{|B_s|<\varepsilon\}}dB_s\Big)^2=\E\int_0^t\tfrac{B_s^2}{\varepsilon^2}\1_{\{|B_s|<\varepsilon\}}ds
\le\int_0^t\P(|B_s|<\varepsilon)\,ds\xrightarrow[\varepsilon\to0]{}0 .
\]

\sub{(c)} Split $\int g_\varepsilon'(B)dB=\int\sgn(B)\1_{\{|B|\ge\varepsilon\}}dB+\int\frac{B}{\varepsilon}\1_{\{|B|<\varepsilon\}}dB$.
The second term tends to $0$ in $L^2$ by (b). The first tends to $\int_0^t\sgn(B_s)dB_s$ in $L^2$, since
$\E\int_0^t\1_{\{|B_s|<\varepsilon\}}ds\to0$. Moreover $0\le g_\varepsilon(x)-|x|=(\varepsilon-|x|)^2/(2\varepsilon)\,\1_{\{|x|<\varepsilon\}}\le\varepsilon/2$.
Hence by (a)
\[
\frac1{2\varepsilon}\Leb\{s\le t:B_s\in(-\varepsilon,\varepsilon)\}=g_\varepsilon(B_t)-g_\varepsilon(B_0)-\int_0^tg_\varepsilon'(B_s)dB_s
\]
converges in $L^2$ to $|B_t|-|B_0|-\int_0^t\sgn(B_s)dB_s$. Thus $L_t$ exists as an $L^2$-limit and
$|B_t|=|B_0|+\int_0^t\sgn(B_s)dB_s+L_t$.

\prob{7.}
\note{The definition should read $T_x=\inf\{t>0:X^x_t=0\}$ with $x>0$, not $\inf\{t:X^x_t>0\}$. In (b) the claim is
$(Y_t)\overset d=(X^1_t)$. All $X^x$, $x>0$, are driven by the same $B$.}

Facts used: $X^x$ is continuous and positive on $[0,T_x)$, and $X^x_t\to0$ as $t\uparrow T_x$ when $T_x<\infty$; also
$B_t/t\to0$ a.s. From the equation we get two more facts.
\begin{itemize}[leftmargin=*,nosep]
\item On $\{T_x=\infty\}$, $X^x$ is unbounded. Indeed, if $X^x_t<R$ for all $t\ge0$ then
$X^x_t\ge x+B_t+at/R\to\infty$, a contradiction.
\item On $\{T_x<\infty\}$, $X^x$ reaches every $r\in(0,x)$.
\end{itemize}
Hence for $0<r<x<R$, $X^x$ does not stay in $(r,R)$ forever.

\sub{(a)} For $t<S=T_x\wedge T_y$, $D_t=X^y_t-X^x_t=y-x+a\int_0^t\big(\frac1{X^y}-\frac1{X^x}\big)ds$ is $C^1$ (the
integrand is continuous) with $D'=a\frac{X^x-X^y}{X^xX^y}=-aD/(X^xX^y)$. Hence
\[D_t=(y-x)\exp\Big(-a\int_0^t\frac{ds}{X^x_sX^y_s}\Big)>0\qquad(t<S).\]
If $T_y<T_x$, then $0<X^x_t<X^y_t\to0$ as $t\uparrow T_y$, so $X^x_{T_y}=0$ by continuity, contradicting $T_y<T_x$.
So $T_x\le T_y$ and $X^x_t<X^y_t$ for $t<T_x$. With $y=x$, the same computation gives $D_t=0\cdot e^{(\dots)}=0$: for
a given path of $B$ the equation has at most one positive continuous solution (pathwise uniqueness), and two solutions
have the same lifetime.

\sub{(b)} Let $W_t=x^{-1}B_{x^2t}$, a BM by scaling. Substituting $s=x^2u$,
$\int_0^{x^2t}ds/X^x_s=x^2\int_0^tdu/X^x_{x^2u}=x\int_0^tdu/Y_u$. Dividing the equation for $X^x_{x^2t}$ by $x$ gives
\[
Y_t=1+W_t+a\int_0^t\frac{du}{Y_u},\qquad t<T_x/x^2 .
\]
So $Y$ is a positive continuous solution of the equation with initial point $1$ driven by $W$, with lifetime
$T_x/x^2$. By pathwise uniqueness the solution with initial point $1$ is a fixed functional of the driving path,
$X^1=\Phi(B)$. Hence $Y=\Phi(W)$ and $Y\overset d=X^1$ (lifetimes included), since $W\overset d=B$. This uses that
$\Phi$ is measurable on path space. That is standard: $Z=X^1-1-B$ solves the pathwise ODE $Z'=a/(1+w+Z)$, whose
solution depends continuously on the path $w$ by Gronwall's inequality. We take it for granted, as the problem does
in asserting that the solution is well defined.

\sub{(c)} Let $s(x)=x^{1-2a}$ for $a\ne\frac12$ and $s(x)=\log x$ for $a=\frac12$. Then $s'(x)=c\,x^{-2a}$ with $c\neq0$,
and $\frac12s''+\frac ax s'=0$. For $0<r<x<R$ let $\tau$ be the exit time of $X^x$ from $(r,R)$; it is a.s.\ finite by
the facts above. On $[0,\tau]$, $X^x$ is an Itô process $dX=\frac aXdt+dB$ with values in $[r,R]$, and $s$ is smooth
there. So Itô's formula gives $s(X_{t\wedge\tau})=s(x)+\int_0^{t\wedge\tau}s'(X)dB$, with integrand bounded by
$\max_{[r,R]}|s'|$, hence a martingale. Bounded convergence ($t\to\infty$) yields $\E s(X_\tau)=s(x)$, i.e.\
\[
p_{r,R}(x):=\P(X^x\text{ hits }r\text{ before }R)=\frac{s(R)-s(x)}{s(R)-s(r)} .
\]
By (a), for $x<y$ we have $T_x\le T_y$ and $X^x<X^y$ on $[0,T_x)$. Properties ii and iv pass from $y$ to every $x<y$,
and properties i and iii pass from $x$ to every $y>x$. So it suffices to prove ii and iv for $x\in\N$, and i and iii for
$x=1/n$, $n\in\N$. These are countably many values, so the almost sure statements hold simultaneously for all $x>0$.

\emph{i.} ($a\ge\frac12$.) If $T_x<\infty$, the path is continuous on $[0,T_x]$, so $\sup_{[0,T_x]}X^x<R$ for some
$R\in\N$. Then $X^x$ hits every $r\in(0,x)$ before $R$. Now $p_{r,R}(x)\to0$ as $r\to0$: for $a=\frac12$,
$s(r)=\log r\to-\infty$; for $a>\frac12$, $s(r)\to+\infty$ while $s(R)-s(x)$ is finite. Hence
$\P(T_x<\infty)\le\sum_{R\in\N,\,R>x}\lim_{r\to0}p_{r,R}(x)=0$.

\emph{ii.} ($a=\frac12$.) By i, $T_x=\infty$. For $0<r<x<R$,
\[\P(X^x\text{ hits }r)\ge p_{r,R}(x)=\frac{\log R-\log x}{\log R-\log r}\xrightarrow[R\to\infty]{}1.\]
Intersecting over $r=1/n$ gives $\inf_{t>0}X^x_t=0$ a.s.

\emph{iii.} ($a>\frac12$.) By i, $T_x=\infty$. Let $M_t=(X^x_t)^{1-2a}>0$ and
$\sigma_n=\inf\{t:X^x_t\notin(\frac1n,n)\}$. Then $\sigma_n\uparrow\infty$, since the path is continuous and positive on
$[0,\infty)$, and $M_{t\wedge\sigma_n}=M_0+\int_0^{t\wedge\sigma_n}s'(X)dB$ is a bounded martingale. Conditional Fatou
gives $\E[M_t|\F_s]\le\liminf_n\E[M_{t\wedge\sigma_n}|\F_s]=\liminf_nM_{s\wedge\sigma_n}=M_s$ (in particular
$\E M_t\le M_0$). So $M$ is a non-negative continuous supermartingale. By Exercise 2, Martingales 5(a) applied to the
submartingale $-M$ (with $\sup_t\E(-M_t)^+=0$), $M_t\to M_\infty\in[0,\infty)$ a.s. Hence $X^x_t\to X_\infty\in(0,\infty]$.
If $X_\infty<\infty$, the continuous path would be bounded on $[0,\infty)$, contradicting the first fact. So
$X^x_t\to\infty$ a.s.

\emph{iv.} ($a<\frac12$.) Here $s(0^+)=0$ and $s(\infty)=\infty$. Fix $R>x$ and let $\tau_r$ ($r<x$) and $\tau_R$ be
the hitting times of $r$ and $R$ ($\inf\emptyset=\infty$). The events $E_r=\{\tau_r<\tau_R\}$ decrease as
$r\downarrow0$. We claim $\bigcap_{r}E_r=\{T_x<\tau_R\}$ up to a null set.

The inclusion $\supseteq$ holds by continuity. Conversely, on $\bigcap_rE_r$ the times $\tau_r$ increase to some
$\sigma\le T_x$, and $\sigma<\infty$ a.s.: if $T_x=\infty$ then $\tau_R<\infty$ by the first fact, and
$\sigma\le\tau_R$. If $\sigma<T_x$, continuity at $\sigma$ would give $X^x_\sigma=\lim_rX^x_{\tau_r}=0$, contradicting
positivity on $[0,T_x)$. So $\sigma=T_x$. Since $\tau_r<\tau_R$ for all $r$, $X^x$ does not reach $R$ before $T_x$,
i.e.\ $T_x<\tau_R$. Hence
\[\P(T_x<\tau_R)=\lim_{r\to0}p_{r,R}(x)=\lim_{r\to0}\frac{R^{1-2a}-x^{1-2a}}{R^{1-2a}-r^{1-2a}}=1-(x/R)^{1-2a}\xrightarrow[R\to\infty]{}1,\]
so $T_x<\infty$ a.s.

\emph{Setting for v and vi.} Fix $y>x>0$ and write $X=X^x$, $D=X^y-X^x$. By (a) and iv, $T_x\le T_y$,
$T_x<\infty$ a.s., and $D$ is $C^1$, positive and non-increasing on $[0,T_x)$. Let $D_{T_x}:=\lim_{t\uparrow T_x}D_t\ge0$.
Since $X^y_t=X_t+D_t\to D_{T_x}$ as $t\uparrow T_x$, we have $T_x=T_y$ iff $D_{T_x}=0$. Moreover, by (a),
\[\log D_{T_x}=\log(y-x)-a\int_0^{T_x}\frac{dt}{X_tX^y_t}.\tag{$\dagger$}\]
On $[0,T_x)$ let $V=X/D>0$, so that $X=DV$ and $X^y=D(V+1)$. Since $1/D$ is $C^1$ with
$d(1/D)=\frac{a}{XX^yD}dt$ (no quadratic covariation with $X$), the product rule gives
\[
dV=\frac1D\Big(\frac aXdt+dB\Big)+X\,d\Big(\frac1D\Big)
=\frac{1}{D^2}\Big(\frac aV+\frac a{V+1}\Big)dt+\frac1D\,dB,\qquad d\langle V\rangle_t=\frac{dt}{D_t^2}.
\]
Hence for $\varphi\in C^2(0,\infty)$, Itô's formula gives
$d\varphi(V)=D^{-2}\big(\tfrac12\varphi''+(\tfrac aV+\tfrac a{V+1})\varphi'\big)(V)\,dt+\varphi'(V)D^{-1}dB$. Let
\[
S(v)=\int_0^v\big(u(1+u)\big)^{-2a}du,\qquad v\ge0,
\]
which is finite since $2a<1$, increasing, and has decreasing derivative $S'(v)=(v(1+v))^{-2a}$. Since
$S''/S'=-2a(\frac1v+\frac1{1+v})$, we get $\frac12S''+(\frac av+\frac a{v+1})S'=0$, so
\[dS(V)=S'(V)D^{-1}dB.\]
For $n\in\N$ let $\rho_n=\inf\{t<T_x:D_t\le\frac1n\text{ or }V_t\le\frac1n\}\wedge n$.

\emph{$\rho_n<T_x$.} If $D_{T_x}<\frac1n$, then $D$ reaches $\frac1n$ strictly before $T_x$. Otherwise
$D\ge\frac1n$ on $[0,T_x)$, so $V\le nX\to0$ and $V$ reaches $\frac1n$ strictly before $T_x$. Thus $V_{\rho_n\wedge\kappa}$
is well defined for every stopping time $\kappa$.

\emph{Optional stopping.} On $[0,\rho_n]$, $|S'(V)D^{-1}|\le S'(\frac1n)\,n$ and $\rho_n\le n$. Hence
$\E S(V_{\rho_n\wedge\kappa})=S(V_0)$ for every stopping time $\kappa$.

\emph{$\rho_n\uparrow T_x$.} $(\rho_n)$ is non-decreasing. If its limit $\rho_\infty$ were $<T_x$, then $D$ and $V$
would be continuous and positive on the compact interval $[0,\rho_\infty]$, hence $\ge c>0$ there. For $n>\max(1/c,\rho_\infty)$
this gives $\rho_n>\rho_\infty$, a contradiction.

\emph{v.} ($\frac14<a<\frac12$.) Now $S(\infty)<\infty$, since $(u(1+u))^{-2a}\le u^{-4a}$ and $4a>1$. Suppose, for
contradiction, that $\P(T_x<T_y)=1$. Then a.s.\ $D_{T_x}>0$ and $X_t\to0$, so $V_t\to0$ as $t\uparrow T_x$. For
$n>\max(1/D_{T_x},T_x,1/V_0)$ we have $D>\frac1n$ on $[0,T_x)$ and $T_x<n$, so $\rho_n$ is the hitting time of $\frac1n$ by $V$,
and $S(V_{\rho_n})=S(\frac1n)\to S(0)=0$. Since $0\le S\le S(\infty)$, bounded convergence gives
$S(V_0)=\lim_n\E S(V_{\rho_n})=0$. This contradicts $V_0=x/(y-x)>0$. Hence $\P(T_x=T_y)>0$.

\emph{vi.} ($a\le\frac14$.) Now $S(\infty)=\infty$, since for $u\ge1$ we have $u(1+u)\le2u^2$, hence $(u(1+u))^{-2a}\ge2^{-2a}u^{-4a}\ge2^{-2a}u^{-1}$ (as $4a\le1$), which is not integrable at $\infty$.

\emph{Step 1: $\sup_{t<T_x}V_t<\infty$ a.s.} For $\lambda>S(V_0)$ let
$\kappa_\lambda=\inf\{t<T_x:S(V_t)\ge\lambda\}$ (with $\kappa_\lambda=T_x$ if this set is empty). On
$\{\kappa_\lambda\le\rho_n\}$ continuity gives $S(V_{\kappa_\lambda})=\lambda$. As $S\ge0$,
$\lambda\P(\kappa_\lambda\le\rho_n)\le\E S(V_{\rho_n\wedge\kappa_\lambda})=S(V_0)$. Since $\rho_n\uparrow T_x$, every
$\omega$ with $\kappa_\lambda<T_x$ satisfies $\kappa_\lambda\le\rho_n$ for large $n$. Hence
$\P(\kappa_\lambda<T_x)\le S(V_0)/\lambda\to0$ as $\lambda\to\infty$. As $S$ is increasing with $S(\infty)=\infty$, this
gives $\sup_{t<T_x}V_t<\infty$ a.s.

\emph{Step 2: $D_{T_x}>0$ on $\{\sup_{t<T_x}V_t<K\}$, for each $K\in\N$ with $K>V_0$.} Let
$\rho'_n=\inf\{t<T_x:D_t\le\frac1n\text{ or }V_t\ge K\}$ (with $\inf\emptyset=T_x$). For $m\ge n$ put
$\theta=\rho'_n\wedge\rho_m$. Then $\theta<T_x$, $\theta\le m$, $D\ge\frac1m$ on $[0,\theta]$, and $0<V\le K$ on
$[0,\theta]$. Integrating the equation for $V$ up to $\theta$, the $dB$-integral has bounded integrand $1/D\le m$ and
bounded time horizon, so it has mean $0$. Dropping the non-negative term $\frac a{V+1}$,
\[
\E\int_0^{\theta}\frac{a}{V_tD_t^2}\,dt\le\E V_\theta-V_0\le K .
\]
As $m\to\infty$, $\theta\uparrow\rho'_n$ (because $\rho_m\uparrow T_x\ge\rho'_n$), so by monotone convergence
$\E\int_0^{\rho'_n}\frac{a\,dt}{V_tD_t^2}\le K$. On $\{\sup_{t<T_x}V<K\}$, $\rho'_n$ is the hitting time of $\frac1n$ by $D$,
which increases to $T_x$ as $n\to\infty$ ($D$ is continuous and positive on $[0,T_x)$). Monotone convergence again
gives $\E\big[\1_{\{\sup V<K\}}\int_0^{T_x}\frac{dt}{V_tD_t^2}\big]\le K/a<\infty$. Since
$\frac{1}{XX^y}=\frac{1}{D^2V(V+1)}\le\frac{1}{D^2V}$, $(\dagger)$ yields $D_{T_x}>0$ a.s.\ on $\{\sup_{t<T_x}V<K\}$.

\emph{Conclusion.} By Step 1, $\bigcup_{K\in\N,\,K>V_0}\{\sup_{t<T_x}V<K\}$ has probability $1$. So $D_{T_x}>0$ a.s., i.e.\
$T_x<T_y$ a.s.

\end{document}
